\section{Related Work and Positioning}
\label{sec:related}

\subsection{Embodied agents and spatial reference}

Embodied conversational agents combine language with a visible body, gesture,
gaze, and shared space \cite{cassell1999embodiment}.  Museum-guide systems show
why the location and social role of an agent matter to interaction
\cite{kopp2005museum,bickmore2011museum}, while recent generative-agent and
multi-agent work adds memory, role specialization, and language-mediated
coordination \cite{park2023generative,li2023camel}.  These systems motivate the
multi-agent setting considered here, but a role or dialogue policy does not by
itself bind an utterance to one Unity entity and one executable operation.

Situated-language research addresses the upstream reference problem.
Approaches combine plan recognition with referring expressions
\cite{gorniak2005grounding}, map directions into perceived environments
\cite{kollar2010directions}, maintain symbolic anchors for discourse
\cite{lemaignan2012grounding}, and learn perceptually grounded word meanings
\cite{kennington2015reference}.  MM-Conv extends this line with synchronized
speech, motion, gaze, and scene context for ambiguous reference in dynamic 3D
dialogue \cite{deichler2026mmconv}.  Our contract does not compete with these
resolvers: it accepts a candidate selection and checks whether the modeled
interface permits that target to reach a specific provider and action.  A
learned or multimodal resolver could therefore replace the current desktop ray
without changing the downstream admission questions.

\subsection{Conversational XR and model-based interface engineering}

Language models increasingly author or operate immersive environments.  LLMR
uses prompting, scene understanding, planning, execution, and self-debugging to
edit mixed-reality worlds \cite{delatorre2024llmr}.  Tell-XR supports
conversational end-user development of event--condition--action automations
\cite{carcangiu2025tellxr}; CUIfy packages LLM-powered conversational agents
for XR \cite{buldu2025cuify}; and AgentHands studies gestures for spatially
grounded agent conversations \cite{liu2026agenthands}.  These systems establish
that conversational generation, agent embodiment, and XR interaction are
valuable and feasible.  The issue studied here is narrower: after a world and
its agents exist, can the interface prevent a selected object, responsible
agent, backend action, and returned evidence from silently diverging?

Model-based interface development has long separated domain, task, abstract
interface, and concrete implementation concerns
\cite{puerta1999framework,calvary2003cameleon}.  UsiXML and MARIA support
multiple development paths and abstraction levels
\cite{limbourg2005usixml,paterno2009maria}, while OpenUIDL addresses runtime
omni-channel interfaces \cite{moldovan2020openuidl}.  Model-driven work for
immersive applications likewise generates AR applications from domain models
\cite{camposlopez2023mdear} or structures VR requirements around roles,
scenes, actions, states, and validation \cite{karre2024vrrequirements}.
\systemname{} is not a general UI description language.  It specializes the
runtime link among spatial entities, embodied-agent responsibility, permitted
capabilities, concrete operations, and observable assertions.

The closest lineage is MLDS.  Fischer et al.\ define a
Meta-Language-defined Structure as a machine-readable artifact governed by a
domain-specific language specification \cite{fischer2025machinereadable}.
Prior work already uses this idea for natural-language-to-Unity materialization
and optional placement refinement \cite{burk2025dynamicvr}, and modular MLDS
instructions add profiles, dependencies, versions, and gatekeeper validation
for evolving toolchains \cite{burk2026modular}.  Consequently, this paper does
not claim novelty for MLDS-to-Unity generation, placement, or modular language
specification.  Its contribution begins at interaction time: the same typed
relations used during materialization remain active as an admission and
evidence contract.

\subsection{Human--AI intelligibility, control, contracts, and provenance}

Human--AI interaction guidelines emphasize communicating system capability and
status, making uncertainty visible, enabling correction, and supporting
recovery \cite{amershi2019guidelines,yang2020reexamining,weisz2024design}.
Lim and Dey organize intelligibility around recurring user questions such as
\emph{What}, \emph{Why}, \emph{Why not}, \emph{What if}, and \emph{How to}
\cite{lim2009demand}; Liao et al.\ adapt this question-driven perspective to
explainable-AI design \cite{liao2020questioning}.  We use these established
questions to organize which contract relation can be exposed at each
interaction gate, rather than introducing a competing explanation taxonomy.
Because explanations do not automatically improve appropriate reliance or
joint performance \cite{bansal2021whole}, the contract supplies inspectable
state while the user study separately evaluates whether the implemented cues
support target, handoff, and responsibility comprehension.

Design by Contract separates preconditions, postconditions, and invariants
\cite{meyer1992contract}; runtime verification evaluates specified properties
against execution observations \cite{leucker2009runtime}.  We adopt this
enforcement intuition without claiming formal program verification.
The relevant obligations cross implementation layers: a request must use a
pinned model, a unique target, a responsible and reachable provider, and an
exact action binding; a response must preserve those identifiers.  W3C PROV
shows how entities, activities, and agents can support exchangeable provenance
\cite{lebo2013provo}, and failure-attribution work demonstrates that output-only
traces are insufficient for multi-agent diagnosis \cite{chen2026elephant}.
Our evidence is intentionally smaller than a full reasoning trace: it records
contract identifiers and observable outcomes, not private chain-of-thought or
an explanation of semantic answer quality.

Requirements traceability is relevant because a link is valuable only when it
supports reasoning about change and failure
\cite{gotel1994traceability,maeder2015benefit}.  The evaluation therefore does
not count links alone.  It tests whether valid cases resolve end to end,
whether controlled mutations are localized, whether runtime requests are
admitted or rejected at the intended boundary, and what representational cost
the links introduce.

\subsection{Positioning}

Table~\ref{tab:positioning} distinguishes the unit constrained by representative
work.  Conversational-XR systems primarily generate or control worlds;
reference-resolution systems infer a target; model-based UI approaches connect
abstract and concrete interfaces; and multi-agent systems coordinate roles.
The present contract starts from a selected candidate and constrains the
specific path that may cross the Unity--backend boundary.

\begin{table*}[t]
  \caption{Positioning by primary unit and assurance locus.  The distinctions
  state the focus of this paper; they do not imply that a cited system lacks
  mechanisms outside its stated contribution.}
  \label{tab:positioning}
  \small
  \begin{tabularx}{\textwidth}{@{}p{27mm}Y Y Y@{}}
    \toprule
    Work & Primary represented unit & Main interface/runtime emphasis &
    Distinction of this work \\
    \midrule
    LLMR \cite{delatorre2024llmr} and Tell-XR
    \cite{carcangiu2025tellxr}
      & World edits or event--condition--action automations
      & Conversational generation and end-user authoring
      & We constrain an already materialized request through target, provider,
        capability, and action. \\
    Spatial grounding \cite{gorniak2005grounding,lemaignan2012grounding}
      & Referring expression, perception, and candidate referent
      & Inferring what a person refers to
      & We begin after candidate selection and test whether the selected entity
        has a permitted executable path. \\
    Model-based UI engineering
    \cite{puerta1999framework,calvary2003cameleon,paterno2009maria}
      & Tasks, domain concepts, and abstract/concrete interfaces
      & Transformation across abstraction and deployment levels
      & We specialize the runtime boundary to spatial entities, agents,
        capabilities, actions, and evidence. \\
    MLDS lineage
    \cite{fischer2025machinereadable,burk2026modular}
      & Machine-readable artifacts and modular language specifications
      & Generation, validation, profiles, and toolchain evolution
      & We keep one typed artifact active during request admission and response
        presentation. \\
    This work
      & Pinned target--provider--capability--action--assertion contract
      & Fail-closed spatial request admission and relation-level evidence
      & The contribution is the composed boundary and its executable,
        corpus-bounded evaluation. \\
    \bottomrule
  \end{tabularx}
\end{table*}
